\setcounter{page}{1}
\section{Introduction}

Architectural floor plans encode a building as a set of labelled rooms. Turning
a scanned or exported plan back into that structured form, one polygon and one
room type per room, is an instance-segmentation problem. It is also the step
that stands between a drawing archive and any later query about areas,
circulation or fire-safety compliance. Published work on the task has
concentrated on residential plans with a handful of coarse categories
\cite{kalervo2019cubicasa5k,kratochvila2024multiunit}.

The dataset used here differs in a way that matters. It consists of German
\emph{Feuerwehrpl\"ane} (fire-brigade plans) of office and institutional
buildings, annotated with 19 functional room categories whose frequencies span
almost three orders of magnitude. Six of those categories (Allg.\
Verkehrsfl\"ache, Abstellraum, Versammlungsraum, Schacht, Garderobe,
Sanit\"arraum) account for 70 of the 1633 polygon annotations, of which 30 fall
in the training split and only 17 in the held-out test split.

That imbalance is the practical obstacle. A detector trained on the raw data
learns the frequent categories and never predicts the rare ones, so the rooms an
inspector most often needs to locate are the ones the model cannot name. The
usual remedy in the long-tailed literature is to reweight the loss or the
sampler \cite{wang2021seesaw}. A more recent remedy is to synthesise the missing
examples with a text-to-image diffusion model
\cite{xie2024mosaicfusion,shin2023fillup,wang2024doge}. Floor plans look like a
favourable setting for the synthetic route, since the drawing is mostly white
space bounded by thin black strokes. A room interior can therefore be repainted
while every wall, door, label and neighbouring room survives untouched. The
project asked one narrow question. Does diffusion or LoRA inpainting of rare
rooms improve room-instance segmentation on this dataset, and does the answer
depend on the backbone?

The work has four parts, all carried out by the author. The first is an
augmentation stack that generates rare-room examples from training-split pixels
only and admits them through automated structural and pixel checks. It produces
fourteen controlled arms covering prompt-based inpainting, an exemplar-guided
variant with copy-paste and feather-blend ablations, and a rank-8 LoRA adapter
\cite{hu2022lora} on Stable Diffusion 1.5 inpainting \cite{rombach2022ldm}. The
second part trains and evaluates every arm twice, once for the 19-class problem
and once for a derived class-agnostic view, using a torchvision Mask R-CNN
\cite{he2017maskrcnn} over three seeds. Those checkpoints were then re-evaluated
under a five-fold cross-validation that pools 21 held-out images instead of 11.
The third part carries the three most promising augmentation datasets over to a
fine-tuned SAM 3 \cite{carion2025sam3}. The fourth separates localisation from
naming: SAM 3 class-agnostic proposals are cropped and classified by a frozen
DINOv2-base encoder \cite{oquab2024dinov2}, with and without the synthetic rare
crops.

The result is negative and consistent. Across three seeds and five folds, no
diffusion or LoRA arm beat the unaugmented baseline by a margin the data
supports. Every paired sign test on the three-seed grid returned $p \ge 0.25$,
and nine of thirteen augmented arms scored below baseline on the class-agnostic
track. The only augmentation with a repeatable positive sign was plain exemplar
copy-paste, the control condition that contains no generated pixels at all. Even
its best case, $+0.045$ class-agnostic AP with all five folds positive, reaches
only $p = 0.0625$, the smallest value attainable at $K = 5$. Rare-class AP
stayed at or near zero everywhere.

What did move the numbers was the backbone. SAM 3 scored $0.713$ class-agnostic
mask AP against $0.624$ for the best Mask R-CNN arm, a gap larger than every
augmentation effect measured in the project. The two-stage experiment shows why
the negative result is still informative. Handed perfect ground-truth crops, the
frozen DINOv2 classifier recalled zero rare instances. The failure therefore
sits in room-type classification at a scale of 17 to 40 test instances. Neither
the segmentation model nor the augmentation pipeline is the limiting factor.
